% ===================================================================
%  CUADRO N.º 14 — La Calle. — Los Comerciantes.
%  Traduction 2026 d'apres texto/io, controlee sur texto/fr et
%  texto/en. Consigne : voir texto/es/10-cuadro-01.tex.
% ===================================================================

%%K t14-apar-1 apar
\VUsaut{4.0mm}
\VUtitre{15.0pt}{60}{CUADRO N.º 14}
\VUsaut{3.0mm}

%%K t14-tit-1 sub
\VUpk{11.4pt}{40}{La Calle. --- Los Comerciantes.}
\VUsaut{3.0mm}

%%K t14-01-1 p
1. --- Mi amigo Juan, que vive en el campo, ha venido a pasar tres días
con mi familia. \VUgras{Hasta} ahora no había tenido nunca ocasión de ver
una gran ciudad, de modo que pasamos todos los días paseando y yo le
explico lo que no conoce.

%%K t14-02-1 p
2. --- Primero fuimos a ver el \VUgras{observatorio} \textsuperscript{(1)}, que le
interesó muchísimo. Al volver por la \VUgras{calle} \textsuperscript{(3)} donde están
los \VUgras{baños turcos} \textsuperscript{(2)}, nos cruzamos con un
\VUgras{entierro} \textsuperscript{(4)}. A la cabeza del \VUgras{cortejo fúnebre}
caminaba el \VUgras{clero}, delante del \VUgras{coche fúnebre} en el que se
había colocado el \VUgras{ataúd}. Muchos amigos seguían a la familia de
\VUgras{luto}.

%%K t14-02-2 p
Cuando hubo pasado el cortejo, cruzamos la calle delante de la gran
\VUgras{cervecería} \textsuperscript{(5)}, donde muchos \VUgras{clientes} bebían
\VUgras{cerveza} en la terraza, al abrigo de la \VUgras{marquesina}
\textsuperscript{(6)} acristalada.

%%K t14-03-1 p
3. --- A Juan le intrigó el \VUgras{escudo de armas} colocado entre la
tienda del \VUgras{comerciante de vinos y licores} \textsuperscript{(7)} y la del
\VUgras{vendedor de música} \textsuperscript{(10)}. Me preguntó qué significaba; le
dije que señalaba el \VUgras{consulado} \textsuperscript{(8)} inglés, y le mostré al
\VUgras{cónsul} \textsuperscript{(9)} de pie en el balcón.

%%K t14-tit-2 sub
\VUpk{10.2pt}{40}{Mirando los escaparates.}
\VUsaut{3.0mm}

%%K t14-04-1 p
4. --- Naturalmente nos detuvimos ante la exposición de
\VUgras{instrumentos de música}, \VUgras{armonios}, \VUgras{órganos} y
\VUgras{pianos}; miramos las cubiertas ilustradas de las \VUgras{partituras}
y las \VUgras{piezas} para piano, y admiramos la \VUgras{lira} de madera
dorada que servía de muestra.

%%K t14-05-1 p
5. --- Las nubes de polvo que levantaban los \VUgras{barrenderos}
\textsuperscript{(11)} y la \VUgras{barredora} \textsuperscript{(12)} nos obligaron a pasar
deprisa por delante de los hermosos \VUgras{juegos de muebles} del
\VUgras{mueblista} \textsuperscript{(13)} y de la exposición de \VUgras{estufas} y
\VUgras{lámparas} del \VUgras{ferretero} \textsuperscript{(14)}.

%%K t14-06-1 p
6. --- En aquel sitio, en efecto, tuvimos que dejar la acera, pues estaba
obstruida por los bultos que se descargaban de un \VUgras{carro de mudanzas}
\textsuperscript{(16)} para llevarlos a una \VUgras{tienda} \textsuperscript{(15)} que acababa
de alquilarse.

%%K t14-06-2 p
Eso nos impidió ver los hermosos coches del \VUgras{carrocero}
\textsuperscript{(17)}, pero no que nos parásemos a mirar al \VUgras{embalador}
\textsuperscript{(19)} que hacía una caja para su vecino, el
\VUgras{vendedor de bicicletas y automóviles} \textsuperscript{(18)}. Después, como
ya era bastante tarde, volvimos a casa.

%%K t14-tit-3 sub
\VUpk{10.2pt}{40}{Un paseo por la ciudad.}
\VUsaut{3.0mm}

%%K t14-07-1 p
7. --- Al día siguiente mi madre propuso que Juan se hiciera un retrato, y
me pidió que lo acompañara a casa del \VUgras{fotógrafo} \textsuperscript{(20)}.
Después del almuerzo fuimos al \VUgras{estudio} del artista, donde tuvimos
que esperar un buen rato. Para pasar el tiempo miramos los
\VUgras{retratos} \textsuperscript{(22)} colgados en la pared.

%%K t14-07-2 p
Cuando nos llegó el turno la cosa no duró mucho, y en cuanto mi amigo
terminó de posar delante de la \VUgras{cámara} \textsuperscript{(21)}, bajamos.

%%K t14-08-1 p
8. --- En el primer piso vimos, por la puerta abierta de la
\VUgras{peluquería} \textsuperscript{(23)}, al oficial que cortaba el pelo a un
cliente.

%%K t14-08-2 p
Ya de nuevo en la calle, pasamos por delante del \VUgras{estanco}
\textsuperscript{(24)}. A Juan le divirtieron tanto el \VUgras{farol} \textsuperscript{(25)}
rojo y la \VUgras{gran pipa} que servía de \VUgras{muestra} \textsuperscript{(26)} que
tropezó con dos ancianos que charlaban tomando un \VUgras{polvo de rapé}
\textsuperscript{(27)}.

%%K t14-tit-4 sub
\VUpk{10.2pt}{40}{La pastelería.}
\VUsaut{3.0mm}

%%K t14-09-1 p
9. --- Los \VUgras{billetes} y las \VUgras{monedas} de todos los países
expuestos en el \VUgras{escaparate} del \VUgras{cambista} \textsuperscript{(29)} nos
retuvieron un momento, pero como era la hora de merendar entramos en la
\VUgras{pastelería} \textsuperscript{(31)} sin detenernos más.

%%K t14-10-1 p
10. --- Ante todas las \VUgras{golosinas} que había en la tienda ---
\VUgras{pasteles, pan de especias, galletas} y \VUgras{dulces} de todas
clases --- apenas podíamos elegir. Al final Juan se decidió por los
\VUgras{pasteles de crema}, y yo tomé sólo una pequeña
\VUgras{tarta de cerezas}, pues los dulces \VUgras{me dan dolor de muelas} y
no tengo ningunas ganas de visitar demasiado a menudo al \VUgras{dentista}
\textsuperscript{(30)}.

%%K t14-tit-5 sub
\VUpk{10.2pt}{40}{La escritura a máquina.}
\VUsaut{3.0mm}

%%K t14-11-1 p
11. --- Para enseñar a mi amigo una \VUgras{máquina de escribir}, de la que
había oído hablar pero que nunca había visto, entramos en el
\VUgras{despacho} \textsuperscript{(32)} de mi \VUgras{primo} \textsuperscript{(34)}. Lo
encontramos dictando sus cartas a su \VUgras{secretaria} \textsuperscript{(35)}, que
es \VUgras{taquígrafa}. Apenas terminada una carta, un
\VUgras{mecanógrafo} \textsuperscript{(33)} la copiaba a máquina. Como no queríamos
interrumpir el trabajo de mi pariente, nos quedamos con él sólo unos
minutos.

%%K t14-12-1 p
12. --- Debajo del despacho de mi primo vive un gran
\VUgras{relojero y joyero} \textsuperscript{(37)} que es además platero. Su
espléndida tienda contiene muchos \VUgras{relojes de pared, relojes de
sobremesa, relojes de bolsillo}, \VUgras{cadenas} y \VUgras{joyas} valiosas,
como \VUgras{collares de perlas}, \VUgras{pulseras} de oro,
\VUgras{pendientes} y \VUgras{sortijas} con \VUgras{piedras preciosas} y
\VUgras{diamantes}. Uno de los escaparates está dedicado por entero a la
\VUgras{platería} y contiene una gran colección de \VUgras{cubiertos} de
plata.

%%K t14-13-1 p
13. --- Al doblar la esquina de la calle, noté que habían cambiado la
\VUgras{placa} \textsuperscript{(36)} que está junto al \VUgras{número} de la casa. Iba
a decirlo en voz alta cuando se oyó el alegre son de una banda militar.
Echamos a correr hacia el regimiento, sin pensar que iba a pasar por
delante de las tiendas del \VUgras{sombrerero} \textsuperscript{(39)} y del
\VUgras{guantero} \textsuperscript{(40)}, y que habríamos estado igual de bien
situados para verlo desfilar quedándonos delante del escaparate del
\VUgras{óptico} \textsuperscript{(38)}, lleno de \VUgras{quevedos, gafas} y
\VUgras{gemelos}.

%%K t14-tit-6 sub
\VUpk{10.2pt}{40}{El desfile.}
\VUsaut{3.0mm}

%%K t14-14-1 p
14. --- A la cabeza del \VUgras{regimiento} \textsuperscript{(41)} marchaba el
\VUgras{tambor mayor} \textsuperscript{(42)}, seguido de los \VUgras{tambores}
\textsuperscript{(43)}, los \VUgras{cornetas} \textsuperscript{(44)} y toda la \VUgras{banda}.
El \VUgras{general} \textsuperscript{(45)}, que estaba en la plaza con su ayudante de
campo, se había detenido junto al \VUgras{surtidor} \textsuperscript{(49)} de la
\VUgras{fuente} \textsuperscript{(48)} para ver el \VUgras{desfile}.

%%K t14-14-2 p
\VUgras{Los oficiales del estado mayor} \textsuperscript{(46)}, el \VUgras{coronel} y
el \VUgras{teniente coronel} lo saludaban con la espada.
\VUgras{Los comandantes} iban a la cabeza de sus \VUgras{batallones}
\textsuperscript{(47)} y cada \VUgras{capitán} mandaba su \VUgras{compañía}, mientras
los \VUgras{tenientes}, los \VUgras{alféreces} y los
\VUgras{suboficiales} ocupaban sus puestos habituales junto a sus hombres.

%%K t14-15-1 p
15. --- Muchos transeúntes se habían reunido en el \VUgras{cruce}
\textsuperscript{(50)}; los comerciantes --- el \VUgras{paragüero} \textsuperscript{(51)}, el
\VUgras{tintorero} \textsuperscript{(53)}, el \VUgras{armero} \textsuperscript{(54)} --- salieron
a ver a los \VUgras{soldados}. Todos los vehículos --- \VUgras{coches de
punto} \textsuperscript{(52)}, \VUgras{coches particulares} \textsuperscript{(57)} y también
la \VUgras{cuba de riego} \textsuperscript{(56)} --- se arrimaron a la acera.

%%K t14-tit-7 sub
\VUpk{10.2pt}{40}{Escenas de la calle.}
\VUsaut{3.0mm}

%%K t14-15-2 p
Un \VUgras{viajante de comercio} \textsuperscript{(55)} que llevaba en la mano sus
\VUgras{maletines de muestras} cruzó la calle a buen paso, y las personas
sentadas en la \VUgras{imperial} \textsuperscript{(59)} del \VUgras{ómnibus}
\textsuperscript{(58)} se volvieron para disfrutar del espectáculo. Un pobre
\VUgras{cantor callejero} \textsuperscript{(60)} con su \VUgras{organillo}
\textsuperscript{(61)} tuvo que interrumpir su \VUgras{canción}, pues el ruido de los
tambores ahogaba su voz.

%%K t14-16-1 p
16. --- Cuando el regimiento llegó delante de la \VUgras{tienda de tejidos}
\textsuperscript{(64)}, las \VUgras{costureras} \textsuperscript{(63)} y los \VUgras{sastres}
\textsuperscript{(62)} dejaron sus \VUgras{máquinas de coser} y su labor para
asomarse a la ventana. \VUgras{El comerciante} \textsuperscript{(65)}, que acababa de
poner \VUgras{trajes} \textsuperscript{(66)} nuevos en su \VUgras{escaparate}, tuvo que
retirar las piezas de \VUgras{paño} \textsuperscript{(67)} y de \VUgras{lienzo}
expuestas en la acera, junto a la \VUgras{boca de alcantarilla}
\textsuperscript{(68)}, por miedo a que la gente las volcara; hasta un viejo
\VUgras{buhonero} \textsuperscript{(69)}, a quien una portera regateaba unas medias,
tuvo que meterse en un portal para acabar su venta.

%%K t14-16-2 p
Acompañamos al regimiento hasta el cuartel \VUgras{y}, muy contentos de
nuestro día, volvimos a casa a la hora de cenar.

%%K t14-tit-8 sub
\VUpk{10.2pt}{40}{Oficios y profesiones diversos.}
\VUsaut{3.0mm}

%%K t14-17-1 p
17. --- Al día siguiente mi padre se ofreció a llevarnos a visitar la
imprenta del periódico local. Aceptamos con entusiasmo.

%%K t14-17-2 p
El \VUgras{impresor} es también \VUgras{librero} \textsuperscript{(70)},
\VUgras{editor}, \VUgras{papelero} y \VUgras{encuadernador}. Ha instalado en
el primer piso la \VUgras{redacción} \textsuperscript{(71)} del periódico, y también
el \VUgras{taller de grabado}, donde trabajan los \VUgras{litógrafos}
\textsuperscript{(72)} y los \VUgras{grabadores en cobre} \textsuperscript{(73)}, y por
último la sala de composición, donde están los \VUgras{cajistas}
\textsuperscript{(74)}. La \VUgras{sala de máquinas} \textsuperscript{(75)} propiamente
dicha, las \VUgras{prensas} y las diversas máquinas, están en la planta
baja.

%%K t14-tit-9 sub
\VUpk{10.2pt}{40}{Juan hace muchísimas preguntas.}
\VUsaut{3.0mm}

%%K t14-18-1 p
18. --- Al salir de la imprenta, Juan se paró, muy intrigado, delante de
una cajita de cristal montada sobre un poste frente a la
\VUgras{mercería} \textsuperscript{(77)}, y preguntó para qué servía. Mi padre le
explicó que era un \VUgras{avisador de incendios} \textsuperscript{(76)}. En efecto,
todo en la ciudad era una maravilla para mi amigo, desde la sencilla
\VUgras{fuente de beber} \textsuperscript{(78)} y el ligero
\VUgras{triciclo de reparto} \textsuperscript{(80)} que montaba un \VUgras{botones}
\textsuperscript{(79)} de librea, hasta el \VUgras{furgón de correos} \textsuperscript{(81)}.

%%K t14-19-1 p
19. --- Mientras mi padre nos dejaba un momento para hablar con el
\VUgras{recaudador de contribuciones} \textsuperscript{(83)}, que se había detenido a
conversar con un \VUgras{abogado} \textsuperscript{(82)} y un \VUgras{notario}
\textsuperscript{(84)}, nos entretuvimos siguiendo el movimiento de la calle. Pasaba
ante nosotros la gente más variada: un \VUgras{aprendiz de pastelero}
\textsuperscript{(85)}, que llevaba en una cesta un espléndido \VUgras{pastel}, venía
detrás de un \VUgras{ropavejero} \textsuperscript{(86)}; un \VUgras{farolero}
\textsuperscript{(87)} hablaba con un \VUgras{gasista} \textsuperscript{(88)}; una
\VUgras{monja} \textsuperscript{(89)} que llevaba de la mano a una huerfanita volvía
a su convento.

%%K t14-tit-10 sub
\VUpk{10.2pt}{40}{De vuelta a casa.}
\VUsaut{3.0mm}

%%K t14-20-1 p
20. --- Justo cuando mi padre volvía con nosotros, Juan soltó una carcajada
al ver las caras tiznadas y la extraña facha de dos
\VUgras{deshollinadores} \textsuperscript{(91)} negros de hollín.

%%K t14-21-1 p
21. --- Yo quería comprar una planta para mi madre en la \VUgras{floristería}
\textsuperscript{(92)}, pero mi padre observó que pesaría mucho para llevarla y que
sería mejor comprar unas \VUgras{naranjas}. Nos acercamos a la
\VUgras{vendedora del puesto} \textsuperscript{(94)}, y cuando el \VUgras{aya}
\textsuperscript{(93)}, que estaba escogiendo algunas para su niño, terminó su
compra, nos sirvieron.

%%K t14-22-1 p
22. --- De vuelta a casa nos encontramos con varios conocidos: una vieja
amiga de mi familia, apoyada en el brazo de su \VUgras{señorita de compañía}
\textsuperscript{(95)}, el \VUgras{ingeniero} \textsuperscript{(96)} de caminos y puentes, y
una de mis primas con su \VUgras{niñera} \textsuperscript{(98)}.

%%K t14-22-2 p
Después pasamos junto a la \VUgras{modista} \textsuperscript{(97)} de mi madre y dos
\VUgras{frailes} \textsuperscript{(99)} forasteros, que se acercaron a preguntar a mi
padre el camino de la estación.
\VUsaut{6.0mm}
\VUfilet{22mm}
